
HESEL modellen (Hot Edge SOL Electrostatic) er en reduceret plasma model udviklet til at beskrive turbulens og transportfænomener i plasmaets rand i en tokamak.
Særligt fokuserer modellen på overgangsområdet mellem det indelukkede plasma og scrape off layer (SOL), hvor magnetfeltlinjerne ikke længere er lukkede.
I dette område transporteres partikler og energi på tværs af magnetfeltlinjerne gennem turbulente strukturer, kaldet blobs. 
HESEL modellen er baseret på drift fluid approksimationen af Braginskii's plasma fluidligninger\cite{udledning_af_hesel}. 

HESEL modellen er reduceret ved at udnytte de karakteristiske tids og længdeskalaer i plasmaets rand, hvilket vedligeholder simulering fysisk realitet og numerisk stabilitet.

I dette afsnit gennemgås de grundlæggende trin i Braginskii's udledning, hvorefter den endelige HESEL model præsenteres.
Formålet er ikke at give en fuldstændig fysisk redegørelse for HESEL modellen, men derimod at introducere de centrale idéer og ligninger, som modellen bygger på.


Udgangspunktet er energi bevarelses sætningen i fase rummet. Denne siger at ændring af energi i faserummet er nul. Da energi og masse er 
raleteret gennem Einsteins berømte ligning $E=mc^2$, kan denne sætning også udtrykkes som ændring af masse $f_a$.
Faserummet består af både partiklernes position $\mathbf{r}$ og hastighed $\mathbf{v}$. For et kollisionsfrit system gælder

\begin{equation*}
    \frac{\mathrm{d}f_a}{\mathrm{d}t}=0.
\end{equation*}
Anvendes kædereglen i faserummet fås:

\begin{equation*}
    \frac{\mathrm{d}f_a}{\mathrm{d}t} = \frac{\partial f_a}{\partial t} + \frac{\mathrm{d}x_\beta}{\mathrm{d}t} \frac{\partial f_a}{\partial x_\beta}
    + \frac{\mathrm{d}v_\beta}{\mathrm{d}t} \frac{\partial f_a}{\partial v_\beta}.
\end{equation*}
Herefter sættes:

\begin{align*}
    \frac{\mathrm{d}x_\beta}{\mathrm{d}t} &= v_\beta, \\
    \frac{\mathrm{d}v_\beta}{\mathrm{d}t} &= \frac{F_{a\beta}}{m_a}
\end{align*}
Når kollisioner inkluderes opnås Boltzmann ligningen

\begin{equation}
    \frac{\partial f_a}{\partial t} + \frac{\partial}{\partial x_\beta} \left(v_\beta f_a\right) +
    \frac{\partial}{\partial v_\beta} \left(\frac{F_{a\beta}}{m_a} f_a\right) = C_a
    \label{eq:boltzmann_plasma}
\end{equation}
Her er $\beta$ indeks for koordinater i det fysiske rum. Det første led, $\partial_t f_a$, beskriver den lokale tidslige ændring af fordelingsfunktionen.
Det andet led, $\partial_{x_\beta}(v_\beta f_a)$, beskriver transport i det fysiske rum.
Det tredje led, $\partial_{v_\beta}\left( \frac{F_{a\beta}}{m_a} f_a \right)$, beskriver transport i hastighedsrummet.
Højresiden $C_a$ beskriver ændringen i fordelingsfunktionen som følge af kollisioner mellem partikler.
Et lille eksempel på anvendelse af \eqref{eq:boltzmann_plasma} kan ses i \eqref{Braginskii example} i appendix B.


I ligningen indgår følgende størrelser:
\begin{align*}
    f_a &= f_a(t,\mathbf{r},\mathbf{v}) &&\text{fordelingsfunktion for partikelart $a$}, \\
    \mathbf{F}_a &= e_a\mathbf{E} + \frac{e_a}{c}\left(\mathbf{v}\times\mathbf{B}\right) &&\text{Lorentzkraften på partikelart $a$}, \\
    m_a &&&\text{massen af partikelart $a$}, \\
    e_a &&&\text{ladningen af partikelart $a$}, \\
    C_a &&&\text{kollisionsleddet for partikelart $a$}.
\end{align*}
Her beskriver $f_a(t,\mathbf{r},\mathbf{v})\mathrm{d}\mathbf{v}\mathrm{d}\mathbf{r}$ antal partikler af art $a$ i det fysiske rum $\mathrm{d}\mathbf{r}$ omkring punktet $\mathbf{r}$.
I resten af udledningen undlades det at specificere hvilken partikelart der er tale om på samme måde Braginskii gør. 
Da kollisioener over alle hasitgheder ikke ændre det totale antal partilkler gælder for kollisions termet.

\begin{equation*}
    \int C\,d\mathbf{v} = 0, 
\end{equation*}
Herefter defineres følgende makroskopiske størrelser:
\begin{align*}
    n(\mathbf{r},t) &= \int f(\mathbf{r},\mathbf{v},t)\,d\mathbf{v}\\
    V_{a}(\mathbf{r},t) &= \frac{1}{n} \int \mathbf{v} f(\mathbf{r},\mathbf{v},t)\,d\mathbf{v} = \langle v \rangle\\
    T(\mathbf{r},t) &= \frac{m}{3}\langle (v - V)^2 \rangle
\end{align*}
Her er $n$, $V$ og $T$ makroskopiske størrelser for tætheden, middelhastigheden og temperaturen for partikler i punktet $\mathbf{r}$ til tid $t$.
For at komme frem til Braginskii's 3 ligninger der beskriver tætheds, impuls og energi ligningerne ganges
\eqref{eq:boltzmann_plasma} med $1$, $m v$ og $\frac{1}{2} m v^2$, integrere over hele hastighedsrummet og
undytte definationer for $n$, $V$ og $T$ fås Braginskii's ligninger.
Selve integrationen og udledningen af ligningerne er ikke udført i denne rapport.

\begin{align}
    &\frac{\partial n}{\partial t} + \nabla \cdot (n \mathbf{V}) = 0 \label{eq:Braginsiis kontinuitetsligning_}\\
    &\frac{\partial}{\partial t}\left(m n \mathbf{V}\right) + \frac{\partial}{\partial x_\beta} \left(m n \langle v v_b \rangle \right) 
    - en\left(E + \frac{1}{c} \left[\mathbf{VB}\right]\right) = \int mv C d\mathbf{v} \label{eq:impulsligning_ikke_simplificeret}\\
    &\frac{\partial}{\partial t} \left(\frac{mn}{2} \langle v^2 \rangle\right)  + \nabla \left(\frac{m n}{2} \langle v^2\mathbf{v} \rangle \right)
    -en\mathbf{EV} = \int \frac{mv^2}{2}Cd\mathbf{v} \label{eq:energiligning_ikke_simplificeret_ikke_simplificeret}
\end{align}
Ligning \eqref{eq:Braginsiis kontinuitetsligning_} kaldes kontinuitetsligningen og beskriver bevarelsen af partikler.
Ligningerne \eqref{eq:impulsligning_ikke_simplificeret} og \eqref{eq:energiligning_ikke_simplificeret_ikke_simplificeret}
kaldes henholdsvis impuls og energiligningen og beskriver bevarelsen af henholdsvis impuls og energi i plasmaet.
\eqref{eq:impulsligning_ikke_simplificeret} og \eqref{eq:energiligning_ikke_simplificeret_ikke_simplificeret} blvier yderligere 
reduceret. 

For hasitgheden af partikeltyper defineres middel hastighed $\mathbf{V}$ og tilfældig hastighed, $\mathbf{v'} =  \mathbf{v - V}$.
Herfra bruger Braginskii idenditetsætningen for to tilfældige variabler, som siger:
\begin{gather*}
    \mathbb{E}\left[XY\right] = \mathbb{E}\left[X\right]\mathbb{E}\left[Y\right] + \mathrm{Cov}(X,Y) \\
    \mathrm{Cov}(X,Y) = \mathbb{E}\left[(X-\mathbb{E}[X])(Y-\mathbb{E}[Y])\right] 
\end{gather*}
Med dette udledes størrelsen
\begin{gather*}
    \left\langle v_\alpha v_\beta \right\rangle = \left\langle(V_\alpha + v_\alpha')(V_\beta + v_\beta')\right\rangle \\
    = V_\alpha V_\beta + V_\alpha \left\langle v_\beta' \right\rangle + V_\beta \left\langle v_\alpha' \right\rangle + \left\langle v_\alpha'v_\beta' \right\rangle
\end{gather*}
Da den tilfældige hastighed er defineret som afvigelsen fra middelhastigheden,
\begin{equation*}
    v_\alpha' = v_\alpha - V_\alpha
\end{equation*}
gælder
\begin{equation*} \left\langle v_\alpha' \right\rangle = \left\langle v_\beta' \right\rangle = 0
\end{equation*}
Dermed fås
\begin{equation*} 
    \left\langle v_\alpha v_\beta \right\rangle = V_\alpha V_\beta + \left\langle v_\alpha'v_\beta' \right\rangle
\end{equation*}
For at simplificere Braginskii's impuls ligning \eqref{eq:impulsligning_ikke_simplificeret} defineres følgene:
\begin{align}
    &\frac{\mathrm{d}f}{\mathrm{d}t}
    = \frac{\partial f}{\partial t} + V_\beta \frac{\partial f}{\partial x_\beta}
    = \frac{\partial f}{\partial t} + (\mathbf{V}\nabla)f \label{op:substantive derivative} \\
    &p = \frac{mn \langle v'^2\rangle}{3} = nT \label{eq:preasure}\\
    &\pi_{\alpha\beta} = nm\langle v_\alpha ' v_\beta ' - \frac{1}{3}v'^2\delta_{\alpha\beta} \rangle \label{eq:stress tensor}\\
    &\mathbf{R} = \int m\mathbf{v}'C \mathrm{d}\mathbf{v} \label{eq:mean change in momentum}
\end{align}
Her er\eqref{op:substantive derivative} operatoren substative afledte anvendt på funktionen $f$,
Ligning \eqref{eq:preasure} definerer det isotrope tryk,
\eqref{eq:stress tensor} den viskøse spændingstensor der  beskriver afvigelser fra isotropt tryk,
\eqref{eq:mean change in momentum} beskriver netto momentumoverførsel som følge af kollisioner mellem partikler.

Braginskii's impuls ligning kan skrives på formen:
\begin{equation}
    mn\frac{\mathrm{d}V_a}{\mathrm{d}t} = - \frac{\partial p}{\partial x_a} - \frac{\partial \pi_{a\beta}}{\partial x_\beta}
    + en\left(E_a + \frac{1}{c} \left[\mathbf{VB}\right]_a\right) + R_a \label{eq:Braginskii's momentum ligning_}
\end{equation}
Venstresiden, af Braginskii's momentum ligning\eqref{eq:Braginskii's momentum ligning_}, repræsenterer ændringen i plasmaets momentum.
Højresiden består af kræfter fra trykgradienter, viskøse spændinger, elektromagnetiske felter og kollisioner med andre partikler.
Disse processer styrer transporten og udviklingen af plasmaets makroskopiske bevægelse.

Braginskii's energi transport ligning \eqref{eq:energiligning_ikke_simplificeret_ikke_simplificeret} er også 
ønsket at skrive på en mere kompakt form. Til dette udledes følgende størrelser:
\begin{align*}
    \langle \frac{v^2}{2} v_\beta \rangle &= \frac{1}{2} V^2 V_\beta + V_\alpha 
    \langle v_\alpha' v_\beta' \rangle + \frac{1}{2} \langle v'^2 \rangle V_\beta +
    \langle \frac{v'^2}{2} v_\beta' \rangle \\
    &= \left( \frac{1}{2} V^2 + \frac{5P}{2mn} \right) V_\beta + 
    \frac{1}{mn} V_\alpha \pi_{\alpha\beta}  + \langle \frac{1}{2} v'^2 v_\beta' \rangle
\end{align*}
Ved at introducere følgende definitioner:
\begin{align*}
    &Q = \int \frac{m}{2} v'^2 C \mathrm{d}\mathbf{v}\\
    &\mathbf{q} = \int \frac{m}{2} v'^2 \mathbf{v}' f \mathrm{d}\mathbf{v} = \langle \frac{1}{2} m v'^2 \mathbf{v}' \rangle \\
\end{align*}
Hvor $Q$ beskriver varme genereret af kollisioner mellem partikler.
$\mathbf{q}$ er flux tætheden for tilfældig transport.
Herfra kan Braginskii's energi ligning skrives på følgende form og endelige form::
\begin{align}
    \frac{\partial}{\partial t}\left(\frac{nm}{2}  V^2 + \frac{3}{2} n T\right) +
    \frac{\partial}{\partial x_\beta} \left( \left(\frac{mn}{2}V^2  +\frac{5}{2} nT\right)V_\beta +
    \left(\pi_{\alpha \beta} V_\alpha \right) + q_\beta \right) = en\mathbf{EV} + \mathbf{RV} + Q \label{eq:Braginskii's energitransport ligning_}
\end{align}
Braginskii's energi transport ligning\eqref{eq:Braginskii's energitransport ligning_} beskriver bevarelsen af 
energi for en plasmaart. Det første led repræsenterer ændringen i den totale energi, som består af summen af
den kinetiske energi og den interne termiske energi. 
Det andet led beskriver energifluxen og indeholder bidrag fra konvektiv energitransport, viskøse spændinger og varmeledning.
Højresiden repræsenterer arbejdet udført af ydre kræfter samt varmeproduktion og energiudveksling forårsaget af kollisioner.

Herved udlede Braginskii sine tre ligninger for tæthed, impuls og energi udvikling for plasma \cite{braginskii1965}. 
\begin{gather}
    \text{Braginskii's kontinuitetsligning}\notag \\
    \frac{\partial n}{\partial t} + \nabla \cdot (n \mathbf{V}) = 0 \label{eq:Braginskii's kontinuitetsligning}\\
    %
    \text{Braginskii's momentumligning}\notag \\
    mn\frac{\mathrm{d}V_a}{\mathrm{d}t} = - \frac{\partial p}{\partial x_a} - \frac{\partial \pi_{a\beta}}{\partial x_\beta}
    + en\left(E_a + \frac{1}{c} \left[\mathbf{VB}\right]_a\right) + R_a \label{eq:Braginskii's momentum ligning}\\
    %
    \text{Braginskii's energitransport ligning}\notag \\
    \frac{\partial}{\partial t}\left(\frac{nm}{2}  V^2 + \frac{3}{2} n T\right) +
    \frac{\partial}{\partial x_\beta} \left( \left(\frac{mn}{2}V^2  +\frac{5}{2} nT\right)V_\beta +
    \left(\pi_{\alpha \beta} V_\alpha \right) + q_\beta \right) = en\mathbf{EV} + \mathbf{RV} + Q \label{eq:Braginskii's energitransport ligning}
\end{gather}

Ved modellering af turbulens i randplasmaer er Braginskii's ligninger imidlertid ofte for komplekse til direkte numerisk anvendelse.
Derfor introduceres den reducerede drift fluid HESEL modeller, hvor en række fysiske antagelser anvendes til at forenkle beskrivelsen af plasmaets dynamik.

I dette projekt anvendes HESEL modellen, som er en to dimensionel drift fluid model udviklet til at beskrive turbulens
og transport i tokamak plasmaers rand og SOL regioner. Modellen bygger på Braginskii's fluidbeskrivelse og fremkommer
gennem en række approksimationer og reduktioner af de oprindelige ligninger.
En detaljeret udledning ligger uden for dette projekts omfang og kan findes i litteraturen \cite{udledning_af_hesel}.
Fokus i dette projekt er derfor ikke på den fulde udledning af HESEL-modellen, men på anvendelsen af modellen
som fysisk grundlag for de efterfølgende numeriske og maskinlæringsbaserede metoder.


HESEL ligningerne er givet på følgende form\cite{udledning_af_hesel}:

\begin{gather}
    \mathrm{d} _t n + n\mathcal{K}(\phi) - \mathcal{K}(p_e) = \Lambda_n  \label{EQ_n} \\ 
    \mathrm{d}_t^0 \omega^* + \left\{\phi, p_i \right\} - \mathcal{K}(p_e+p_i)=\Lambda_{\omega^*} \label{EQ_w} \\
    \frac{3}{2}\mathrm{d} _t p_e + \frac{5}{2}p_e \mathcal{K}(\phi) -\frac{5}{2}\mathcal{K}\left(\frac{p_e^2}{n}\right)= \Lambda_{p_e}\label{EQ_p_e} \\
    \frac{3}{2}\mathrm{d} _t p_i + \frac{5}{2}p_i \mathcal{K}(\phi) +\frac{5}{2}\mathcal{K}\left(\frac{p_i^2}{n}\right) - p_i\mathcal{K}(p_e+p_i) = \Lambda_{p_i} \label{EQ_p_i}
\end{gather}


\section{Lambda termer}
Lambda termerne er givet ved:
\begin{align}
\Lambda_n =& D_e\nabla_\perp^2 n -\frac{n}{\tau} -\frac{n-n_p}{\tau_p} \notag\\ 
            &-\alpha\left(\tilde T_e+ \frac{\bar T_e}{\bar n}\tilde n-\tilde\phi\right) \label{L_n}\\ 
\Lambda_\omega^{*}=& D_i\nabla_\perp^2\omega^{*}-\frac{\omega^{*}}{\tau}+\frac{\rho_s}{L_c}\left[1-\exp\left(\phi_m-\frac{\phi_s}{T_{e,s}}\right)\right] \notag\\
            &-\alpha\left( \tilde T_e+ \frac{\bar T_e}{\bar n}\tilde n -\tilde\phi \right) \label{L_w}\\
\Lambda_{p_e} =& \frac{5}{2}D_e\nabla_\perp^2 p_e  + \left(\frac{16}{6} - \frac{5}{2}\right)\nabla\cdot\left(n\nabla_\perp T_e\right) - \frac{9p_e}{2\tau} - \frac{T_e}{\tau^{SH}}-\Theta -\frac{p_e-p_{e,p}}{\tau_p} \notag\\
            &-\bar T_e \alpha\left( \tilde T_e + \frac{\bar T_e}{\bar n}\tilde n -\tilde\phi \right)\label{L_p_e}\\
\Lambda_{p_i} =& D_i\nabla_\perp^2 p_i - D_iT_i\nabla_\perp^2n - \frac{9p_i}{2\tau} +\Theta -\frac{p_i-p_{i,p}}{\tau_p} +p_i\Lambda_{\omega^*} \label{L_p_i}
\end{align}

\section{Funktioner}
Funktionerne er gevet ved:
\begin{align*}
    &n: n(x,y,t): &\text{Tæthed} \\
    &p_{e,i}: p_{e,i}(x,y,t): &\text{Elektron- og iontryk}  \\
    &T_{e,i}: T_{e,i}(x,y,t) = p_{e,i}/n: \quad &\text{Elektron- og iontemperatur} \\
    &\phi: \phi(x,y,t): \quad &\text{Elektrisk potentiale} \\
    &\omega^*: \omega^*(x,y,t) = \nabla^2\phi + \nabla^2 p_i: \quad &\text{Vorticitet} \\
\end{align*}


\section{Operatorer}
Herved er operatorerne defineret med:
\begin{gather}
    \mathcal{K}(f)= - \frac{\rho_s}{R}\partial_y f \label{curvature}\\
    \left\{ f, g\right\} = \partial_x f \partial_y g - \partial_y f \partial_x g \label{poission}\\
    \nabla = \left(\partial_x, \partial_y, \partial_z \right)^T \label{nabla}\\
    \nabla_\perp = \left(\partial_x, \partial_y, 0 \right)^T \label{nable_perp}\\
    \nabla^2 = \left(\partial^2_x, \partial^2_y, \partial^2_z \right)^T\label{nabla2}\\
    \nabla_\perp^2 = \left(\partial^2_x, \partial^2_y, 0 \right)^T\label{nabla2_perp}\\
    \mathrm{d}_t^0 f = \partial_t f + \left\{\phi, f \right\} \label{dt0}\\
    \mathrm{d}_t f = \partial_t f  + \frac{1}{B}  \left\{ \phi , f \right\} \label{dt} \\    
    \bar f = \frac{1}{L_y}\int_0^{L_y} f \mathrm{d}y \label{bar}\\
    \tilde f = f - \bar f\label{tilde}\\
\end{gather}


\section{Koeffecienter}
Samling af koeffecienter:
\begin{gather}
    \Theta = \frac{3m_e}{m_i}v_{ei}\left(p_e - p_i\right)\label{Theta}\\
    \alpha = \frac{2T_{e0}}{v_{ei}m_eL^2_\parallel\omega_{i0}}\label{alpha}\\
    \phi_m = \ln(\sqrt{m_i/2\pi m_e})\label{phi_m}\\
\end{gather}


\section{Profile region}
Hvor er profile regionen defineret ved:
\begin{gather}
    \frac{n-n_p}{\tau_p}\sigma_{lcfs} \label{profile_n}\\
    \frac{p_e-p_{e,p}}{\tau_p}\sigma_{lcfs} \label{profile_p_e}\\
    \frac{p_i-p_{i,p}}{\tau_p}\sigma_{lcfs} \label{profile_p_i}\\
    \frac{\phi_s}{T_{e,s}} = \frac{\phi(x,y,t)}{T_e(x,y,t)}*\sigma_{connected} \label{phi_s_connected}\\
    \frac{\phi_s}{T_{e,s}} = \frac{\bar\phi(x,y,t)}{\bar{T_e}(x,y,t)}*\sigma_{disconnected} \label{phi_s_disconnected}\\
\end{gather}


\newpage
